\section{Spectral Definition of Mass as a Persistence Measure of a Route \(1\) Coupled to a Physical Regime}
Mass is the spectral measure of the persistence of a route when a genealogical section is coupled to a physical regime. It is not a scalar attribute affixed in advance to a particle name. The APP--TRIT--TPK dynamics constructs the route, its memory, its signature, and its character; Dimensional Mechanics realizes that genealogy on a resonance mass line and determines, through the coupling, the observable subspace and its spectrum.

An elementary particle is represented by a persistent section; a noncentral family, by a stable multisection; a composite particle, by a composition of genealogical registers with a linking boundary; a resonance, by a pole of the coupled operator; a topological sector, by its fusion ring and its braiding representation before any mass chart; a virtual section, by its terminal horizon of prolongation. These realizations are distinct and must not be flattened into a single list of numbers.

The external magnitude never enters the selection of cell, route, signature, reader, tower coefficient, or eigenvector. Metrological comparison is performed only after the internal output has been constructed and fixed.

\Needspace{7\baselineskip}
\section{Complete state space of mass}
Let \(\mathcal C_{81}\) be the cellular atlas. Each cell admits four bidirectional route orientations, so the finite space of oriented routes is

\[
 \mathcal R_{324}
 =\{(c,h_0,v_0):c\in\mathcal C_{81},\ (h_0,v_0)\in\{\pm1\}^2\}.
\]
The cellular projection, the family relation, and the coarse stratification produce

\[
 |\mathcal C_{81}|=81,\qquad
 |\mathcal F_{56}|=56,\qquad
 |\mathcal M_{13}|=13,\qquad
 |\mathcal R_{324}|=324.
\]
The electron occupies the unique central fixed point. The twelve noncentral multisections are finite fibers stable under cellular inversion and do not, in general, admit an equivariant point section. Their natural object is therefore an operator on the fiber rather than a number chosen from among its cells. The null sector is incorporated through the zero section of the mass line, not as the value of a positive character.

The physical output catalog distinguishes three charged leptons, three neutrino interaction realizations, six quark flavors with their color orbits, two null gauge realizations, two massive gauge bosons, one scalar sector, two composite families, three topological codomains, and a system of virtual sections. This subsequent classification does not limit the number of composite or resonant states: composition and excitation generate additional families within their respective codomains.

\Needspace{7\baselineskip}
\section{Genealogical trajectory of a mass realization}
The law does not begin with a particle label nor with a number. Its order is

\[
\boxed{
\begin{aligned}
\mathrm{APP}\to\mathrm{TRIT}\to\mathrm{TPK}
&\to\widetilde\Gamma^{\rm surv}
\to\Gamma^{\rm obs}
\to L_{108}
\to\sigma\in\mathbb Z^5\\
&\to\chi_0
\to\chi_{\rm full}
\to\widehat{\mathcal M}
\to\text{coupling}\\
&\to\text{spectrum/pole}
\to\mathcal L_M
\to\text{comparison}.
\end{aligned}}
\]
Here:

\begin{itemize}
\item APP arranges the local incompatibility, the sum/product sheets,
and their orientation;
\item the TRIT preserves sign, threshold, and bidirectional opening;
\item the TPK evolves phase, memory, carry, survival, and return;
\item a particle is a persistent section of an extension, and its observable route is
the history it leaves in the genealogical record;

\item the integer signature is additive;
\item the character turns additive composition into a positive ratio;
\item the towers refine the character through a single global harmonic operator;
\item the mass line supplies the dimensional type;
\item measurement is coupling: the coupling determines which subspace and which spectrum are
physically accessible, without consulting the quantity to be measured subsequently.

\end{itemize}
Mass does not select the route. The persistent route, the genealogical record, and the coupling select the spectral object whose mass is read.

\Needspace{7\baselineskip}
\subsection{Conserved components of the enriched state of the TPK along the mass trajectory}
The state that feeds the mass law is

\[
 X_K=(w_6,w_{30},R_{36},g_K,B_K,p_K,c_K,\varepsilon_K,\omega_K),
\]
where short word, prolonged word, region, nonadic phase, block memory, prefix, carry, sheet, and winding are preserved. The chain

\[
 w_6\longrightarrow w_{30}\longrightarrow R_{36}
 \longrightarrow G_9
\]
It does not directly project a mass: it constructs the persistent identity that is later read.

\Needspace{7\baselineskip}
\subsection{Nonadic filtration with memory}
From each state one obtains a refinement fiber formed by \(729\) subcylinders,

\[
 \{C_{K,j}:0\le j<729\}.
\]
The pruning \(G_9\) satisfies

\[
 G_9(C_{K,j})\in\{0,1\},\qquad
 \sum_{j=0}^{728}G_9(C_{K,j})=1.
\]
If \(j_K\) is the survivor,

\[
 P_{K+1}=729P_K+j_K.
\]
After nine steps,

\[
 g(K+9)=g(K),
\]
but they change, in general,

\[
 (B_{K+9},P_{K+9},\operatorname{pref}_{K+9},
 \operatorname{Witt}_{K+9})
 \ne
 (B_K,P_K,\operatorname{pref}_K,\operatorname{Witt}_K).
\]
This is the equational form of persistence: the phase returns, not the state. The mass can recognize a particle after the return because the genealogical record preserves the difference.

\Needspace{7\baselineskip}
\subsection{Nonadic carry transduction to the signed dodecaphasic state}
For a temporal coordinate \(\theta\), define

\[
 g_0(\theta)=\theta\bmod9,\qquad
 \beta(\theta)=
 \left\lfloor\frac{\widetilde\theta}{9}\right\rfloor\bmod12.
\]
Then

\[
 g_0(\theta+9)=g_0(\theta),\qquad
 \beta(\theta+9)=\beta(\theta)+1\pmod{12}.
\]
Thus, the nonadic return preserves the phase and advances one memory tooth. The 108 times are

\[
 108=9\cdot12,
\]
so that each genealogical record contains a complete cycle of both scales. The 36.864 carry compositions describe the local combinations of this bidirectional memory.

\Needspace{7\baselineskip}
\subsection{Completion and Conservation of the Unit}
The base boundary mil decomposes as

\[
 999=729+270=729+243+27,
\]
while the completion is

\[
 1000=729+271=729+243+27+1,
\qquad
 271=270+1.
\]
\(+1\) is the unique vertex lying entirely on the boundary and realizes the carry that returns the unit. The partition \(729/271\) distinguishes the active kernel from the mediating boundary; it is not a golden-ratio approximation. The same conservation allows a route to change blocks without losing its normalized identity.

\Needspace{7\baselineskip}
\subsection{Theorem of the complete oriented atlas}
Every cell \(c\in\mathcal C_{81}\) has four orientations

\[
 (h_0,v_0)\in\{\pm1\}^2.
\]
The map

\[
 (c,h_0,v_0)\longmapsto\Gamma(c,h_0,v_0)
\]
is bijective over \(\mathcal R_{324}\). Consequently,

\[
 |\mathcal R_{324}|=4|\mathcal C_{81}|=324.
\]
Each route has 108 genealogical register entries, so the total temporal census is

\[
 81\cdot108=8748.
\]
The quadruple

\[
 (\text{seed row},\text{seed column},h_0,v_0)
\]
determines of manner single the route, its word \(w_6\), its word dodecaphasic, their returns, memory, carry, signature historical and readers. Therefore, the 324 routes not only may is enumerated: may is reconstructed and is joined exactly, field by field, without using a mass.

\Needspace{7\baselineskip}
\subsection{Quotients of families and multisections}
The quotient of the 81 cells by route history yields 56 families. Their multiplicities are distributed as

\[
 41\times1,\quad10\times2,\quad2\times3,\quad1\times4,\quad2\times5,
\]
and

\[
 41+20+6+4+10=81.
\]
The stratification into thirteen multisects possesses ranges

\[
 1,\quad2,\quad4,\quad6,\quad8,\quad12,\quad16
\]
with multiplicities

\[
 1,\quad2,\quad3,\quad2,\quad3,\quad1,\quad1,
\]
and sum of ranks

\[
 1+4+12+12+24+12+16=81.
\]
The electron occupies the rank-one block. The remaining sectors lie in higher-rank blocks and require a representation projector, not a nominal cell.

\Needspace{7\baselineskip}
\subsection{Signature, readers, and memory by route}
For each \(\Gamma\in\mathcal R_{324}\), the following are calculated:

\[
 \sigma_\Gamma\in\mathbb Z^5,\qquad
 K_D(\Gamma),\qquad K_\Omega(\Gamma),
\]
together with first returns of class, orbit, cell, and kinematic state; separation of prefixes; carry debt; ambiguity; stabilized triads; sheet and winding. The family \(K_D\) has 14 profiles; \(K_\Omega\), nine; together they form 40 pairs and 18 coincidences. This diversity is the fine information that makes it possible to pass from a coarse fiber to a physical operator.

\Needspace{7\baselineskip}
\section{From the APP defect to the full signature}
\Needspace{7\baselineskip}
\subsection{Primordial defect and unfolding}
At the center APP,

\[
 B_c=9P_++5P_-,\qquad 9-5=4.
\]
The primordial local gap is 4. The global defect \(54\) belongs to a later unfolding \(2\to6\to54\); it must not be projected back onto the center as though it were the same operation.

\Needspace{7\baselineskip}
\subsection{\(12\) eventos and divisor dodecafasico}
The oriented coordinate does not receive a factor \(1/6\) arbitrarily. The chain is

\[
12\ \text{events}
\to s_C^{\rm tot}
\to 1+5+6
\to s_C^{(12)}=\frac{s_C^{\rm tot}}6
\to W_\pm=6n_A\pm s_C
\to\operatorname{SNF}(1,12)
\to\chi_{\rm mass}.
\]
The \(1/6\) is the oriented quotient of the closure of twelve events. It is not the residue (-1/12), it is not the pseudoinverse of a valency, and it is not the hexad--octad incidence.

\Needspace{7\baselineskip}
\subsection{Construction of the complete full signature from the route register}
The genealogical record CT--108 produces

\[
 \sigma(\Gamma)
 =\bigl(n_A,s_C,k_{\Delta_4},b_{120},b_\zeta\bigr)\in\mathbb Z^5.
\]
The coordinates have distinct origins:

\begin{itemize}
\item \(n_A\): dominant duration/persistence;
\item \(s_C\): net orientation of the bilayer, divided by six after the
dodecaphasic closure;

\item \(k_{\Delta_4}\): minimal-torsion channel;
\item \(b_{120}\): event coefficient of the register;
\item \(b_\zeta\): step-three autocorrelation.
\end{itemize}
We define

\[
 \zeta_{270}:=135\Delta_4^2,
\]
and remains separated from

\[
 s_{270}=q_-^{270}-q_+^{270}.
\]
The 324 oriented routes are joined bijectively to their five historical coordinates through the canonical route identity. This join uses only genealogical structure and consults no metrological magnitude.

\Needspace{7\baselineskip}
\subsection{Direct angle, counter-angle, and conjugate channels}
The direct angular coordinate proceeds from the already constructed fine structure constant:

\[
 A=1000\alpha_{\rm HMT}
 =7.297352569283800997285105472380663\ldots{}^\circ.
\]
The counter-angle is not obtained from a mass or from the Barbero--Immirzi functional. Let

\[
 D_A=\exp\!\left(-\frac{100\pi}{9}\alpha_{\rm HMT}\right),
\]
\[
 \mathcal C_\pi(\alpha_{\rm HMT})
 =169\alpha_{\rm HMT}^{6}
 +\frac{D_A\alpha_{\rm HMT}^{7}}
        {1-D_A\alpha_{\rm HMT}},
\]
\[
 y_*=\frac{\pi}{90}(H_5+\alpha_{\rm HMT})
      -\mathcal C_\pi(\alpha_{\rm HMT}),
 \qquad
 C^*=\frac{180}{\pi}y_*.
\]
Therefore,

\[
 C^*=2.123738338968645724261951380393\ldots{}^\circ.
\]
The coordinate \(A\) is the common angular mode and \(C^*\), the oriented differential mode. Its passage to the Archimedean chart is

\[
 x=A_{\rm rad}=\frac{\pi A}{180},
 \qquad
 y=C^*_{\rm rad}=\frac{\pi C^*}{180}.
\]
In the real algebra generated by the involution \(R=\operatorname{diag}(1,-1)\), with \(P_\pm=(I\pm R)/2\), the biangular block is

\[
 B_{A,C^*}=AI+C^*R
 =(A+C^*)P_++(A-C^*)P_-.
\]
Angle plus counter-angle and angle minus counter-angle are thus the two eigenvalues of a single object. Exponentiation produces the channels

\[
 q_+=e^{-x-y},\qquad q_-=e^{-x+y}.
\]
The chart is invertible:

\[
 x=-\frac12\log(q_+q_-),
 \qquad
 y=\frac12\log\frac{q_-}{q_+}.
\]
The transformation \(C^*\mapsto-C^*\) preserves the projectors and exchanges the channels. The return coordinate satisfies, after the regional construction,

\[
 \eta_{\rm ret}=\frac{C^*}{2}-\alpha_{\rm HMT}\,{\rm deg},
\]
without constituting a second selection rule for \(C^*\).

\Needspace{7\baselineskip}
\subsection{Central Character}
With

\[
 \Theta_0=
 \left(A_{\rm rad},\frac{C^*_{\rm rad}}6,
 \Delta_4,s_{120},\zeta_{270}\right),
\]
the central character is

\[
 \chi_0(\sigma)=
 \exp\!\left(
 n_AA_{\rm rad}+\frac{s_C}{6}C^*_{\rm rad}
 +k_{\Delta_4}\Delta_4+b_{120}s_{120}+b_\zeta\zeta_{270}
 \right).
\]
By additivity of the signature,

\[
 \chi_0(\sigma_1+\sigma_2)=\chi_0(\sigma_1)\chi_0(\sigma_2).
\]
The physical unit does not enter this exponential: \(\chi_0\) is a positive dimensionless automorphism of the mass line.

\Needspace{7\baselineskip}
\section{Global harmonic hierarchy of towers}
With the coordinates (x=\allowbreak{}A\_\{\textbackslash{}rm rad\}) and (y=\allowbreak{}C\textasciicircum{}*\_\{\textbackslash{}rm rad\}) already constructed, let

\[
 q_\pm=\exp\!\left[-\frac{\pi}{180}(A\pm C^*)\right]
 =e^{-x\mp y}.
\]
In a globally consistent convention,

\[
 T=e^{-xI-yR},\qquad
 c_n=\operatorname{Tr}(T^n)=q_-^n+q_+^n,
\]
\[
 s_n=-\operatorname{Tr}(RT^n)=q_-^n-q_+^n.
\]
Equivalently and operatorially,

\[
 c_n=2e^{-nx}\cosh(ny),\qquad
 s_n=2e^{-nx}\sinh(ny).
\]
The involution \(C^*\mapsto-C^*\) fixes the projectors and permutes the spectral coefficients, or equivalently the associated channels. It follows that

\[
 c_{m+n}=\frac{c_mc_n+s_ms_n}{2},\qquad
 s_{m+n}=\frac{s_mc_n+c_ms_n}{2},
\]
\[
 c_n^2-s_n^2=4e^{-2nx},
\]
and the common second-order recurrence. The admissible heights form

\[
 \mathfrak S=\langle90,120\rangle
 =\{90,120\}\cup\{30r:r\ge6\}.
\]
The values \(90,120,180,210,240,270,360,420\) are genealogical witnesses, not the complete tower. Before the quotient, height 360 preserves two histories:

\[
 k[\mathfrak S_0]
 \simeq k[X_{90},Y_{120}]/(X_{90}^4-Y_{120}^3).
\]
The general elevation is written

\[
 \chi_{\rm full}(\sigma,\eta)
 =\chi_0(\sigma)
  \exp\!\left(\sum_{h\in\mathfrak S}\eta_hs_h\right),
\]
with

\[
 N_{120}=b_{120}+\eta_{120}.
\]
The numerical sum does not erase the dual provenance: \(b_{120}\) comes from the ledger and \(\eta_{120}\) from the spectral refinement. Nor can it absorb \(b_\zeta\zeta_{270}\) into \(\eta_{270}s_{270}\).

\Needspace{7\baselineskip}
\section{Absolute Rama: action, clock and mass line}
The determinantal reader produces

\[
 \operatorname{SNF}(H_4)=\operatorname{diag}(1,2,2,4),\qquad
 |\operatorname{coker}H_4|=16,
\]
\[
 \Sigma_H=\varphi\alpha_{\rm HMT}^{16},\qquad
 \operatorname{ord}_{10}(\Sigma_H)=-34.
\]
The two oriented sections of a single action line are

\[
 \hbar_{\rm pre}=H_5\,10^{-34}\mathcal U_S,
 \qquad
 \hbar_{\rm ret}=\eta_{\rm ret}\,10^{-34}\mathcal U_S.
\]
Its additive fine chart and its multiplicative chart are

\[
 \delta_{2w}=H_5-\eta_{\rm ret},\qquad
 R_{\rm act}=\frac{H_5}{\eta_{\rm ret}}>0.
\]
They are not two Planck constants. The common decade fixes the absolute order; the quotient transports relative outward/return information.

Declared \(t_0\in\mathcal L_T^+\), CT--108 fixes

\[
 \omega_0=\frac{2\pi}{108t_0},\qquad
 \varepsilon_{\rm clk}=\hbar_{\rm ret}\omega_0,
 \qquad
 m_{\rm clk}=\varepsilon_{\rm clk}c_{\rm int}^{-2}.
\]
The character acts by homothety on \(\mathcal L_M\); it does not become \(\hbar\). In mass ratios, every common factor cancels:

\[
 \frac{m_Y}{m_X}
 =\chi_{\rm full}(\sigma_Y-\sigma_X,\eta_Y-\eta_X)
 R_{\rm act}^{\kappa_Y-\kappa_X}.
\]
Thus \(10^{-34}\) fixes the absolute scale but does not correct a relative residual. A relative modification can follow only from the route, the specific clock, the reader, or the towers.

Remains a distinction necessary. HMT constructs a section internal a time declared \(t_0\); the coordinate MeV requires a chart. For close the chart absolute without a convention externa must construirse

\[
 \Gamma\longmapsto t_\Gamma
 \quad\text{or}\quad
 \Gamma\longmapsto\omega_\Gamma.
\]
Section 10 defines a clock reader based on returns and holonomy. Its absolute normalization is determined when the route operator fixes a positive section of the temporal line.

\Needspace{7\baselineskip}
